%%%%%%%%%%%%%%%%%%%%%%%%%%%%%%%%%%%%%%%%%%%%%%%%%%%%%%%%%%%%%%%%%%%%%%%%%%%%%%%%
%%%%%%%%%%%%%%%%%%%%%%%% ------------------------------ %%%%%%%%%%%%%%%%%%%%%%%%
%%%%%%%%%%%%%%%%%%%%%%%% | Formato escuela - INF PUCV | %%%%%%%%%%%%%%%%%%%%%%%%
%%%%%%%%%%%%%%%%%%%%%%%% | by Sebastian Garcia        | %%%%%%%%%%%%%%%%%%%%%%%%
%%%%%%%%%%%%%%%%%%%%%%%% | @sebaignacioo              | %%%%%%%%%%%%%%%%%%%%%%%%
%%%%%%%%%%%%%%%%%%%%%%%% |            2023            | %%%%%%%%%%%%%%%%%%%%%%%%
%%%%%%%%%%%%%%%%%%%%%%%% ------------------------------ %%%%%%%%%%%%%%%%%%%%%%%%
%%%%%%%%%%%%%%%%%%%%%%%%%%%%%%%%%%%%%%%%%%%%%%%%%%%%%%%%%%%%%%%%%%%%%%%%%%%%%%%%
%
%  TRABAJO DE INVESTIGACIÓN 2026 · ICI-5444
%  Empresa 6 · TERABYTE · Tema TI-05 · Serverless y computación en el borde
%
%  ARCHIVO ÚNICO. Todo el informe vive en este main.tex, en el orden
%  definitivo acordado por el equipo (ver Terabyte-Guia-Equipo.md).
%  Cada bloque lleva en un comentario: responsable, presupuesto de
%  extensión y recordatorios de las Indicaciones (FEP00.3.26).
%
%  Buscar "%% >>>" para saltar de bloque en bloque.
%  Buscar "TODO" para ver lo que falta.
%
%  Preámbulo, clase y estilos sin cambios respecto de la plantilla.
%  La versión anterior (un .tex por sección) quedó en _legacy/.
%
%%%%%%%%%%%%%%%%%%%%%%%%%%%%%%%%%%%%%%%%%%%%%%%%%%%%%%%%%%%%%%%%%%%%%%%%%%%%%%%%

\documentclass[12pt, letterpaper, table, xcdraw]{inf-pucv}

\usepackage{import}
\usepackage{styles}
\usepackage{subcaption}
\usepackage{pdflscape}   % \begin{landscape} para el cuadro comparativo (3.2)

%% \anexo{X}{Título}: abre un anexo (o documento adjunto) con prefijo X.
%% Reinicia la numeración de páginas (X1, X2, ...) y la de tablas y figuras
%% (Tabla X.1, Figura X.1) para que no hereden la del capítulo 5.
\newcommand{\anexo}[2]{%
    \clearpage
    \setcounter{page}{1}%
    \renewcommand{\thepage}{\hspace*{1em}#1\arabic{page}}%
    \setcounter{table}{0}\setcounter{figure}{0}%
    \renewcommand{\thetable}{#1.\arabic{table}}%
    \renewcommand{\thefigure}{#1.\arabic{figure}}%
    \phantomsection
    \addcontentsline{toc}{section}{#2}%
    \section*{#2}%
}

\renewcommand{\cftdotsep}{4}
\renewcommand{\cftfigpresnum}{\bfseries Figura }
\cftsetindents{figure}{0em}{5.5em}
\renewcommand{\cfttabpresnum}{\bfseries Tabla }
\cftsetindents{table}{0em}{5.5em}
\setlength{\parindent}{1cm}

\begin{document}

%%%%%%%%%%%%%%%%%%%%%%%%%%%%%%%%%%%%%%%%%%%%%%%%%%%%%%%%%%%%%%%%%%%%%%%%%%%%%%%%
%% >>> ÍNDICE DE BLOQUES (orden definitivo)
%%
%%   Portada e índice general
%%   Resumen ejecutivo ........................... Daniel
%%   Introducción y contexto ..................... Daniel   (cierra con párrafo de aporte)
%%   1  Marco teórico y conceptual ............... Matías
%%   2  Desarrollo y análisis técnico
%%      2.1 Modelos de ejecución y aislamiento ... Matías
%%      2.2 Arranque en frío ..................... Isidora
%%      2.3 Facturación y punto de equilibrio .... Francisca
%%      2.4 Diseño orientado a eventos ........... Valentina
%%      2.5 WebAssembly y WASI ................... Isidora
%%      2.6 Límites y dependencia del proveedor .. Fabián
%%   3  Análisis comparativo de tecnologías y productos
%%      3.1 Alternativas adicionales ............. Claudio y Daniel
%%      3.2 Criterios y cuadro comparativo ....... Valentina (criterios), Rodolfo (AWS/Azure,
%%                                                  consolida), Claudio y Daniel (resto)
%%   4  Caso de estudio e implementación práctica
%%      4.1 Prueba de concepto ................... Vicente
%%      4.2 Modelo de costos ..................... Francisca
%%      4.3 Caso de licitación TERABYTE .......... Fabián
%%   5  Tendencias y riesgos ..................... Fabián
%%   Conclusiones y recomendación ................ todo el grupo
%%   Bibliografía ................................ todos (>= 10 académicas + técnicas)
%%   Anexo A  Declaración de uso de IA ........... Daniel recopila; cada uno declara lo suyo
%%   Anexo B  Código y mediciones de la PoC ...... Vicente
%%   Anexo C  Planilla del modelo de costos ...... Francisca
%%   Anexo D  Búsqueda documentada de alternativas Claudio y Daniel
%%   Documento adjunto: cuestionario ............. todos (3-4 preguntas c/u); Daniel arma índice
%%
%% Presupuesto total del cuerpo (resumen -> conclusiones): 10 a 15 páginas.
%% Referencia: 1 página de texto corrido ~ 450 palabras.
%%%%%%%%%%%%%%%%%%%%%%%%%%%%%%%%%%%%%%%%%%%%%%%%%%%%%%%%%%%%%%%%%%%%%%%%%%%%%%%%


%%%%%%%%%%%%%%%%%%%%%%%%%%%%%%%%%%%%%%%%%%%%%%%%%%%%%%%%%%%%%%%%%%%%%%%%%%%%%%%%
%% >>> PORTADA
%%%%%%%%%%%%%%%%%%%%%%%%%%%%%%%%%%%%%%%%%%%%%%%%%%%%%%%%%%%%%%%%%%%%%%%%%%%%%%%%

\thispagestyle{empty}
\pagenumbering{gobble}

\begin{center}
    \membrete{%
        PONTIFICIA UNIVERSIDAD CATÓLICA DE VALPARAÍSO\\
        FACULTAD DE INGENIERÍA\\
        ESCUELA DE INGENIERÍA INFORMÁTICA\\
    }

    \vspace*{3.5cm}

    \tituloPortada{%
        Serverless y computación \\
        en el borde
    }

    \vspace*{2.5cm}

    \alumnos{%
        Francisca Antonia Abarca Pezo\\
        Vicente André Arratia Caroca\\
        Isidora Valentina Cisternas Cisternas\\
        Rodolfo Antonio Fernández Vera\\
        Valentina Ignacia Guzmán Elgueta\\
        Daniel Ignacio Miranda Vázquez\\
        Matias Ignacio Reyes Pizarro\\
        Fabián Andres Solís Piñones\\
        Claudio Patricio Toledo Mac-lean\\
    }

    \vspace*{1.4cm}

    \datosPortada{%
        \datoPortada{Empresa}{N.\textordmasculine\ 6 --- TERABYTE}
        \datoPortada{Tema}{TI-05}
        \datoPortada{Asignatura}{Taller de Formulación de Proyectos Informáticos}
        \datoPortada{Fecha de entrega}{21 de septiembre de 2026}
    }

    \vspace*{1.2cm}

    \fechaPortada{Septiembre de 2026}
\end{center}

\clearpage

%% >>> ÍNDICE GENERAL
\tableofcontents
\clearpage

%% Listas de figuras y tablas: PENDIENTE de decisión del equipo.
%% Descomentar si se decide incluirlas.
% \phantomsection
% \addcontentsline{toc}{section}{Lista de figuras}
% \listoffigures
% \clearpage
% \phantomsection
% \addcontentsline{toc}{section}{Lista de tablas}
% \listoftables
% \clearpage


%%%%%%%%%%%%%%%%%%%%%%%%%%%%%%%%%%%%%%%%%%%%%%%%%%%%%%%%%%%%%%%%%%%%%%%%%%%%%%%%
%% >>> RESUMEN EJECUTIVO                                  Responsable: Daniel
%% Presupuesto: 0,5 página / ~200 palabras. Se escribe AL FINAL.
%% Contenido: qué se comparó, resultado principal (volumen de equilibrio,
%% latencia medida), recomendación en una frase. Sin contexto ni definiciones.
%% IA: solo corrección de estilo (nivel 1), declarada en Anexo A.
%%%%%%%%%%%%%%%%%%%%%%%%%%%%%%%%%%%%%%%%%%%%%%%%%%%%%%%%%%%%%%%%%%%%%%%%%%%%%%%%

\pagenumbering{arabic}
\setcounter{page}{1}

\phantomsection
\addcontentsline{toc}{section}{Resumen ejecutivo}
\section*{Resumen ejecutivo}

% TODO (Daniel)

\clearpage


%%%%%%%%%%%%%%%%%%%%%%%%%%%%%%%%%%%%%%%%%%%%%%%%%%%%%%%%%%%%%%%%%%%%%%%%%%%%%%%%
%% >>> INTRODUCCIÓN Y CONTEXTO                            Responsable: Daniel
%% Presupuesto: 0,75 página / ~350 palabras. Se escribe AL FINAL.
%% Contenido: el problema (serverless/edge vs. always-on en costo, latencia y
%% complejidad operacional), el proyecto de licitación de TERABYTE como caso,
%% cómo se organiza el informe. Sin historia del cloud.
%%
%% PÁRRAFO FINAL OBLIGATORIO (punto 4 Indicaciones): distinguir qué viene de
%% la ficha TI-05 y qué es aporte del grupo. Enumerar: alternativas
%% adicionales (3.1), PoC con mediciones propias (4.1), aplicación al caso
%% TERABYTE (4.3), riesgos y evidencia contradictoria (5), recomendación
%% comprometida. El profesor usa este párrafo para evaluar el análisis crítico.
%% IA: prohibida en el párrafo de aporte (punto 6.1).
%%%%%%%%%%%%%%%%%%%%%%%%%%%%%%%%%%%%%%%%%%%%%%%%%%%%%%%%%%%%%%%%%%%%%%%%%%%%%%%%

\phantomsection
\addcontentsline{toc}{section}{Introducción y contexto}
\section*{Introducción y contexto}

% TODO (Daniel): cuerpo de la introducción.

% TODO (Daniel): párrafo de identificación del aporte (va aquí, al cierre).

\clearpage


%%%%%%%%%%%%%%%%%%%%%%%%%%%%%%%%%%%%%%%%%%%%%%%%%%%%%%%%%%%%%%%%%%%%%%%%%%%%%%%%
%% >>> 1. MARCO TEÓRICO Y CONCEPTUAL                      Responsable: Matías
%% Presupuesto: máximo 1 página / ~450 palabras.
%% Contenido: definiciones mínimas para leer el capítulo 2: serverless (qué es
%% y qué no es), FaaS, contenedores sin servidor, edge computing, instancia
%% siempre encendida, y las tres dimensiones de comparación (costo, latencia,
%% complejidad operacional). Fuentes: CNCF, papers, docs oficiales.
%% No explicar qué es un contenedor, una VM o la nube.
%% Referencias cruzadas a 2.x van con número fijo porque esas subsecciones
%% aún no tienen \label. Si se les agrega label, cambiar por \ref.
%%
%% IA: borrador revisado con apoyo de IA a partir de fuentes abiertas y
%% verificadas en el original (ver TRAZABILIDAD al final).
%%%%%%%%%%%%%%%%%%%%%%%%%%%%%%%%%%%%%%%%%%%%%%%%%%%%%%%%%%%%%%%%%%%%%%%%%%%%%%%%

\section{Marco teórico y conceptual}
\label{sec:marco}

\textbf{Computación sin servidor.} Es la ejecución de aplicaciones sin gestionar servidores: una plataforma ejecuta, escala y factura el código según la demanda de cada momento. No implica que desaparezcan los servidores ni la operación, sino que su gestión pasa al proveedor \parencite{cncf2018}. \textcite{schleiersmith2021} la caracterizan por tres cualidades: abstracción de los servidores y de su operación, cobro por uso en vez de por reserva (sin cargo por recursos ociosos) y escalamiento automático desde cero. Tampoco es sinónimo de funciones: incluye los servicios de \textit{backend} gestionados, o BaaS \parencite{golec2024}.

\textbf{Modelos de ejecución.} En las \textit{funciones como servicio} (FaaS) la unidad es una función breve, activada por un evento o una solicitud HTTP, que el proveedor escala horizontalmente \parencite{cncf2018,leitner2019} y libera tras un periodo sin uso, lo que se llama escalamiento a cero \parencite{golec2024}. En los \textit{contenedores sin servidor} la unidad es una imagen de contenedor y el equipo ya no aprovisiona ni escala máquinas virtuales \parencite{awsfargate}, aunque no siempre hay escalamiento a cero: Cloud Run retira la última instancia cuando no hay solicitudes \parencite{gcprun}, mientras que Fargate cobra hasta que la tarea termina \parencite{awsfargatepricing}. El \textit{cómputo en el borde} procesa en recursos ubicados entre la fuente de los datos y el centro de datos de la nube \parencite{shi2016}; este informe estudia su variante sin servidor \parencite{aslanpour2021}.

\textbf{Modelo de referencia.} Las \textit{instancias siempre encendidas} son capacidad aprovisionada en forma continua (máquinas virtuales, o contenedores sobre nodos fijos) que se factura por tiempo encendido, se use o no \parencite{ustiugov2021}. Con cargas en ráfagas, esa capacidad queda subutilizada \parencite{eismann2021}.

\textbf{Dimensiones de comparación.} El \textit{costo} en \textit{serverless} se calcula por consumo (invocaciones, tiempo de ejecución, memoria, CPU y transferencia de datos) \parencite{ghorbian2026}, de modo que el tiempo ocioso no se paga \parencite{adzic2017}; en el modelo siempre encendido es casi fijo para una capacidad dada. Esa diferencia define un volumen de equilibrio (secciones~2.3 y~\ref{sec:costos}). La \textit{latencia} es el tiempo entre el envío de una solicitud y su respuesta, suma de la espera en cola, la asignación de recursos y la ejecución \parencite{golec2024}. Al escalar a cero se agrega el \textit{arranque en frío}, el retardo de preparar un entorno nuevo \parencite{ustiugov2021}, tratado en las secciones~2.2 y~\ref{sec:poc}; el borde, en cambio, reduce la latencia de red al acercar el cómputo a la fuente de los datos \parencite{golec2024}. La \textit{complejidad operacional} es el trabajo del equipo para desplegar, escalar, monitorear y mantener el sistema. El modelo sin servidor traslada parte de ese trabajo al proveedor, un motivo frecuente de adopción \parencite{eismann2021}, pero \textcite{leitner2019} documentan herramientas de prueba y despliegue poco maduras y la necesidad de componer servicios del proveedor, a lo que se suman sus límites y la dependencia (secciones~2.4 y~2.6). Cómo se mide cada dimensión se justifica en la sección~\ref{sec:cuadro}.

%% ---------------------------------------------------------------------------
%% TRAZABILIDAD (qué respalda cada cita):
%%  cncf2018 ........ definición textual de serverless; "does not mean that we
%%                    no longer use servers"; FaaS "triggered by events or
%%                    HTTP requests"; BaaS.
%%  schleiersmith2021 tres "essential qualities" (abstracción, pay-as-you-go vs.
%%                    reservation-based, escalamiento desde cero).
%%  golec2024 ....... (PDF del proyecto) serverless integra BaaS y FaaS (§1);
%%                    scale to zero y cold start (§1, §3.3); definición de
%%                    latencia como cola + asignación + ejecución (§3.2);
%%                    edge acerca el cómputo a la fuente (§1).
%%  leitner2019 ..... funciones "atomic and short-running" escaladas
%%                    horizontalmente; herramientas de prueba/despliegue poco
%%                    maduras; FaaS como "glue" de servicios (abstract).
%%  awsfargate ...... "no longer have to provision, configure, or scale
%%                    clusters of virtual machines".
%%  gcprun .......... "even the last remaining instance will be removed".
%%  awsfargatepricing cobro desde "download your container image" hasta que
%%                    "the task terminates" (no se cita ningún precio).
%%  shi2016 ......... definición de "edge" (§II).
%%  aslanpour2021 ... serverless edge computing (abstract).
%%  ustiugov2021 .... (PDF del proyecto) VMs "billed for their up-time
%%                    regardless of usage"; cold start tras inactividad (§1, §2.1).
%%  eismann2021 ..... ahorro en cargas en ráfagas con baja utilización en
%%                    hosting tradicional; menos carga operacional (deployment,
%%                    scaling, monitoring) como motivo de adopción.
%%  ghorbian2026 .... (PDF del proyecto) métricas de precio: execution time,
%%                    memory, request count, CPU, bandwidth (§2.3).
%%  adzic2017 ....... "application idle time is free".
%% ---------------------------------------------------------------------------

\clearpage


%%%%%%%%%%%%%%%%%%%%%%%%%%%%%%%%%%%%%%%%%%%%%%%%%%%%%%%%%%%%%%%%%%%%%%%%%%%%%%%%
%% >>> 2. DESARROLLO Y ANÁLISIS TÉCNICO
%% Seis subtemas obligatorios de la ficha TI-05.
%% Presupuesto: 0,75 página / 300-350 palabras POR SUBSECCIÓN.
%% Estructura sugerida de cada subsección (tres párrafos):
%%   1) el concepto, con fuente;
%%   2) cómo se comporta en cada modelo (FaaS / contenedor serverless / edge /
%%      always-on), con fuente;
%%   3) "cuándo conviene": la conclusión comparativa del subtema.
%% Nada de descripciones de productos, capturas ni listas de características.
%% Los precios NO se buscan aquí: se citan desde la Tabla 3.2.
%%%%%%%%%%%%%%%%%%%%%%%%%%%%%%%%%%%%%%%%%%%%%%%%%%%%%%%%%%%%%%%%%%%%%%%%%%%%%%%%

\section{Desarrollo y análisis técnico}

%%%%%%%%%%%%%%%%%%%%%%%%%%%%%%%%%%%%%%%%%%%%%%%%%%%%%%%%%%%%%%%%%%%%%%%%%%%%%%%%
%% >>> 2.1 MODELOS DE EJECUCIÓN Y AISLAMIENTO                BORRADOR
%%
%% Extensión: párrafos 1 y 2 = ~320 palabras. Con el párrafo 3 (50-70
%% palabras, autoría humana) queda en ~380, un poco sobre el presupuesto de
%% 300-350. Si falta espacio, lo primero que podemos sacar sin perder el hilo
%% es la frase de Cloud Run functions (\parencite{gcprunfunctions}).
%%
%% Referencias cruzadas a 2.2 y 2.5 van con número fijo porque esas
%% subsecciones aún no tienen \label (mismo criterio del marco teórico).
%%%%%%%%%%%%%%%%%%%%%%%%%%%%%%%%%%%%%%%%%%%%%%%%%%%%%%%%%%%%%%%%%%%%%%%%%%%%%%%%

\subsection{Modelos de ejecución y aislamiento}
\label{sec:modelos}

Además de la unidad que ejecuta (sección~\ref{sec:marco}), cada modelo se distingue por cómo aísla a los clientes que comparten una máquina \parencite{shafiei2022}. Compartir servidores permite vender tramos cortos de capacidad que de otro modo quedaría ociosa \parencite{liuniu2024}, pero multiplica el costo de toda sobrecarga de aislamiento \parencite{agache2020}. \textcite{agache2020} agrupan los mecanismos en tres familias: los \textit{contenedores}, que comparten el núcleo del anfitrión y obligan a transar entre seguridad y compatibilidad; la \textit{virtualización}, que da a cada carga su propio núcleo y aísla con fuerza, pero ocupa más memoria y tarda segundos en arrancar \parencite{wen2023}; y el \textit{aislamiento por lenguaje}, como los \textit{isolates} de V8, que comparten un proceso a costa de compatibilidad y de exposición a ataques de canal lateral. En medio se ubican las microVM, máquinas virtuales mínimas, y gVisor, un núcleo en espacio de usuario que atiende las llamadas al sistema del contenedor \parencite{wang2022}.

En FaaS, AWS Lambda crea cada entorno en una microVM de Firecracker que no reutiliza entre funciones ni entre cuentas \parencite{awslambdasec}; según sus creadores, Firecracker agrega menos de 5~MB por microVM, arranca en menos de 125~ms, ejecuta binarios Linux sin modificar y sostiene también a Fargate \parencite{agache2020}. Cloud Run, base de las Cloud Run functions \parencite{gcprunfunctions}, ofrece un entorno basado en gVisor, de arranque rápido pero que no emula todas las llamadas al sistema, y otro basado en microVM, totalmente compatible con Linux y con mejor CPU y red, pero de arranque más lento en algunos servicios \parencite{gcprunenv}; ambos suman una capa de virtualización por hardware \parencite{gcprunsec}. En pruebas independientes, gVisor aisló mejor el rendimiento, pero penalizó más las llamadas al sistema y la red \parencite{wang2022}. En el borde, con nodos de pocos recursos donde VM y contenedores resultan pesados \parencite{xie2021}, Cloudflare Workers ejecuta miles de \textit{isolates} por máquina, solo acepta JavaScript y WebAssembly (sección~2.5) y, contra Spectre, restringe relojes e hilos \parencite{cfworkerssec}. En el modelo siempre encendido, la frontera es la VM arrendada: compatibilidad total \parencite{agache2020} y recursos reservados aunque estén ociosos \parencite{liuniu2024}.

Ningún mecanismo aventaja a los demás en todos los frentes a la vez, ya que seguridad, compatibilidad, densidad y velocidad de arranque se transan entre sí, tal como concluyen tanto \textcite{agache2020} como \textcite{wang2022}, quienes afirman textualmente que «ningún runtime es ventajoso en todos los aspectos». Google mismo reconoce que su entorno de segunda generación (microVM) «generalmente rinde más bajo carga sostenida», pero arranca más lento que el basado en gVisor en ciertos servicios \parencite{gcprunenv}, lo que entrega un criterio directo ligado a la condición de carga. Los \textit{isolates} de V8 solo admiten JavaScript y WebAssembly, sin binarios nativos \parencite{cfworkerssec}, mientras que Firecracker sí ejecuta binarios Linux sin modificar \parencite{agache2020}, lo que marca la restricción de compatibilidad de código. \textcite{wang2022} también muestran que gVisor penaliza especialmente las llamadas al sistema, la red y los archivos pequeños. Y aunque Firecracker arranca en menos de 125~ms, Lambda igual mantiene un \textit{pool} de microVM prearrancadas, porque ese tiempo no alcanza para el escalamiento en ráfaga. La elección entre contenedores, microVM, isolates o VM completa depende, entonces, del tipo de código que se ejecutará (lenguaje, binarios, llamadas al sistema), del patrón de carga (ráfagas cortas versus carga sostenida) y de si el costo de mantener recursos siempre listos compensa evitar la latencia de arranque.

%% ---------------------------------------------------------------------------
%% OPCIONAL: Claude me tiro hacer una tabla resumen (1/4 de página). que nos seria util para la presentación y
%% para el cuestionario, incluirla en el informe solo si cabe en 15 páginas.
%% todas las cifras salen de las fuentes citadas en cada fila.
%% ---------------------------------------------------------------------------
% \tabla{Mecanismos de aislamiento en plataformas sin servidor y de borde}{aislamiento}{|p{2.6cm}|p{3.3cm}|p{4.4cm}|p{3.6cm}|}{%
%     \filaTabla{\textbf{Mecanismo} & \textbf{Frontera entre clientes} & \textbf{Ventaja y costo} & \textbf{Dónde se usa}}
%     \filaTabla{microVM (Firecracker) & Núcleo propio bajo un monitor de VM mínimo & $<$5~MB de sobrecarga, $<$125~ms de arranque, binarios Linux sin cambios \parencite{agache2020} & Lambda, Fargate; Cloud Run 2.\textsuperscript{a} gen. \parencite{gcprunenv}}
%     \filaTabla{Núcleo de aplicación (gVisor) & Llamadas al sistema atendidas en espacio de usuario & Arranque rápido; no emula todas las llamadas; sobrecarga en llamadas al sistema y red \parencite{gcprunenv,wang2022} & Cloud Run 1.\textsuperscript{a} gen.}
%     \filaTabla{Isolate V8 & Memoria separada dentro de un proceso compartido & Miles por máquina; solo JavaScript y WebAssembly \parencite{cfworkerssec} & Cloudflare Workers}
%     \filaTabla{VM completa & Máquina virtual arrendada & Compatibilidad total; recursos reservados aunque estén ociosos \parencite{agache2020,liuniu2024} & Modelo siempre encendido}
% }

%% ---------------------------------------------------------------------------
%% 2.2                                                    Responsable: Isidora
%% Causas (carga del entorno, init del runtime, aislamiento), medición,
%% mitigaciones (concurrencia aprovisionada, instantáneas/SnapStart, entornos
%% livianos, min instances). Cierra citando los resultados propios:
%% "las mediciones del grupo (sección 4.1) ..." -> \ref{sec:poc}.
%% ---------------------------------------------------------------------------
\subsection{Arranque en frío: causas, medición y mitigaciones}
\label{subsec:arranque_en_frio}

El escalado a cero libera los entornos inactivos; según mediciones de 2021, AWS desaloja la mitad de sus contenedores cada 380~s~\cite{copik2021sebs}. Sin instancias disponibles, la plataforma debe reservar recursos, crear el entorno aislado (ver~2.1), cargar runtime y dependencias y ejecutar la inicialización~\cite{golec2025coldstart}: esa latencia es el arranque en frío; reutilizar una instancia preparada es un arranque en caliente. Crece con dependencias grandes y runtimes pesados como la JVM o .NET, frente a Python o Node.js~\cite{golec2025coldstart}, y aparece también al escalar de $N$ a $N+1$ instancias en ráfagas~\cite{giannini2022cpuboost}.

Medirlo exige separar inicialización y ejecución comparando invocaciones frías y calientes, como hace Invoc-Overhead~\cite{copik2021sebs}. AWS Lambda registra esa Init Duration y, desde agosto de 2025, la factura en todas las configuraciones: el arranque en frío también es costo~\cite{gupta2025initbilling}.

Las mitigaciones actúan sobre distintas partes del proceso y difieren en costo. La concurrencia aprovisionada lo evita con instancias ya inicializadas (AWS, Cloud Run, Azure Functions)~\cite{aws2023provisionedconcurrency, googlecloudrunmininstances, azurefunctionsflex}, pero paga capacidad ociosa, acercándose al modelo siempre encendido (ver~2.3 y~3.2), y no cubre el paso de $N$ a $N+1$, que Google aborda con startup CPU boost~\cite{giannini2022cpuboost}. Las instantáneas, como SnapStart, se saltan la inicialización restaurando una captura cifrada de la microVM~\cite{aws2022snapstart}, aunque parte del retraso pasa a la ejecución al cargarse la memoria desde disco~\cite{ustiugov2021}. Los entornos livianos (V8, WebAssembly; ver~2.5) abaratan la creación del entorno: en un prototipo, Wasm redujo el arranque en frío cerca de un 94\,\% frente a Docker~\cite{gackstatter2022pushing}, pero exige adaptar el código.

El impacto depende del patrón de carga~\cite{copik2021sebs}. En APIs síncronas, donde degrada la experiencia de usuario, sirven las instantáneas (con costo de caché en Python y .NET)~\cite{aws2022snapstart} o los aislamientos livianos, cambiando de plataforma~\cite{gackstatter2022pushing}; aun así, Cloudflare afirmó en 2020 haberlo eliminado~\cite{cloudflare2020workers}, pero reapareció con scripts más pesados y su técnica de sharding lo redujo 10 veces~\cite{cloudflare2025shard}. En tráfico sostenido, la concurrencia aprovisionada es más económica sobre el 60\,\% de uso continuo~\cite{gupta2025initbilling}. En procesos asíncronos, fuera de la ruta crítica, la latencia es tolerable: es despreciable en funciones de más de 10~s y, sin alta utilización, escalar a cero cuesta menos que una máquina virtual~\cite{copik2021sebs}.%% ---------------------------------------------------------------------------
%% 2.3                                                    Responsable: Francisca
%% GB-segundo, vCPU-segundo, tiempo de CPU vs. tiempo de reloj, número de
%% invocaciones, salida de datos (egress); concepto de punto de equilibrio
%% frente a instancias reservadas / contenedores siempre encendidos.
%% El cálculo numérico va en 4.2 (\ref{sec:costos}); aquí los conceptos y
%% variables. Toda cifra que aparezca se toma de la Tabla 3.2 con su fecha.
%% ---------------------------------------------------------------------------
\subsection{Facturación por consumo y punto de equilibrio}
\label{sec:facturacion}

Según \textcite{lin_demystifying_2026}, la factura de una plataforma sin servidor combina cuatro factores: el tiempo de reloj facturable, los recursos asignados o consumidos durante ese tiempo, la granularidad o el mínimo de cobro y un cargo fijo por invocación. Las funciones cobran invocaciones y duración ponderada por la memoria asignada; AWS Lambda redondea al milisegundo y, desde agosto de 2025, factura también la fase de inicialización (\cite{Shubham_aws_2025}), mientras que Azure Functions ofrece un plan Premium con instancias precalentadas cobrado por núcleo-segundo y memoria (\cite{microsoft2026functionsscale}). Los contenedores sin servidor cobran vCPU y memoria por segundo: Fargate con un mínimo de un minuto por tarea (\cite{aws_precios2026}) y Cloud Run por solicitud o por instancia, a elección (\cite{google2026cloudrunpricing}). En el borde, Cloudflare Workers cobra solicitudes y tiempo de CPU consumido, no de reloj, y no cobra transferencia de datos (\cite{aws_cloudflare_2026}). Los valores numéricos se remiten en las Tablas 3.2a y 3.2b (sección~\ref{sec:cuadro}).

Como cada modelo mide una unidad distinta, la CPU facturada es entre 1,01 veces (Workers) y 3,63 veces (Google) la consumida en simulaciones sobre trazas de Huawei Cloud (\cite{lin_demystifying_2026}). El cobro por reloj incluye la espera por E/S (por ejemplo, esperar a que responda una base de datos), y encadenar funciones duplica esa facturación \cite{eismann2021}. En AWS Lambda, Google e IBM incluyen también la inicialización  (\cite{lin_demystifying_2026}). Workers cobra la CPU consumida sin contar la duración, con un tope de 128 megabytes por {\itshape Isolate} (\cite{cloudflare2026workerslimits}). Por eso las unidades solo se comparan normalizadas (Tablas 3.2a y 3.2b).

El punto de equilibrio es la carga desde la cual pagar por consumo cuesta más que una instancia siempre encendida. Pagar por consumo conviene mientras la prima de precio sea menor que la razón entre tiempo contratado y tiempo usado (\cite{xiayue_charles_lin_serverless_2020,copik2021sebs}). El cruce llega antes con carga regular que en ráfagas, y todos los patrones cruzan el punto de equilibrio si el volumen crece (\cite{liu2024demystifying}). La capacidad fija solo conviene si se mantiene una alta utilización (\cite{copik2021sebs}). Sin variar el volumen, el cruce se desplaza con los descuentos por reserva (37,5\% a un año y 59\% a tres años en EC2, derivados de la Tabla 3.2a), la concurrencia aprovisionada (\cite{lin_demystifying_2026}) y costos como mantenimiento y egreso (\cite{hamza2024costdynamics}). La evidencia es frágil: dos casos anecdóticos (\cite{gojko_adzic_serverless_2017}), precios de 2020 (\cite{xiayue_charles_lin_serverless_2020}) (\cite{copik2021sebs}) y trazas de un solo proveedor (\cite{lin_demystifying_2026}). El detalle de su cuantificación está en la sección~\ref{sec:costos}.
%% ---------------------------------------------------------------------------
%% 2.4                                                    Responsable: Valentina
%% Colas, publicación/suscripción, idempotencia, entrega al menos una vez,
%% orquestación frente a coreografía. Aquí también se discute la complejidad
%% operacional (observabilidad, depuración distribuida). Candidatos "otros"
%% de datos serverless -> pasar a Claudio y Daniel (3.1).
%% ---------------------------------------------------------------------------
\subsection{Diseño orientado a eventos}

Una arquitectura orientada a eventos construye sistemas débilmente acoplados que
colaboran emitiendo y respondiendo a eventos en lugar de invocarse entre sí
\parencite{awseventbridge}. CloudEvents define el evento como un registro de
datos que expresa una ocurrencia y su contexto, enrutado hacia los consumidores
interesados sin identificar un destino \parencite{cloudevents102}. Para
\textcite{eugster2003}, la fortaleza del estilo está en el desacoplamiento
simultáneo en espacio, tiempo y sincronización entre publicador y consumidor.

La cola punto a punto entrega cada mensaje a un único consumidor, con semántica
\emph{uno-de-n}, y al extraerse de forma síncrona no desacopla en sincronización
\parencite{eugster2003}. En publicación/suscripción, suscribirse a un tópico
equivale a integrar un grupo y publicar, a difundirlo entre sus miembros
\parencite{eugster2003}. El patrón fija quién recibe el evento; la semántica de
entrega, cuántas veces llega y qué debe tolerar el consumidor. Las dos variantes
de SQS lo muestran: la estándar entrega al menos una vez y obliga a manejar
duplicados con operaciones idempotentes, la FIFO procesa exactamente una vez
\parencite{awssqsqueuetypes}. El RFC 9110 llama idempotente al método cuyo efecto
pretendido sobre el servidor no cambia al repetir la petición
\parencite{rfc9110}. En las fuentes revisadas el exactamente-una-vez nunca
aparece sin condiciones: Kafka lo alcanza sumando al productor idempotente el
envío transaccional a varias particiones \parencite{kafkadesign}; Pub/Sub, solo
dentro de una región y en suscripciones \emph{pull}
\parencite{gcppubsubexactlyonce}.

La coordinación admite dos formas opuestas: la orquestación centraliza el orden
de los pasos y la coreografía deja que cada servicio reaccione por su cuenta a
los eventos ajenos. La guía prescriptiva de AWS (documentación técnica de
terceros, declarada como tal) recomienda coreografía dentro de un microservicio y
orquestación al cruzar ese límite, con el patrón \emph{saga} para revertir
transacciones distribuidas \parencite{awscoordination}. El criterio es el costo
del error: MQTT reserva QoS~0 a sensores ambientales, donde perder una lectura no
importa, y QoS~2 a facturación, donde un duplicado es un cobro incorrecto
\parencite{mqttv5}; Step Functions limita el exactamente-una-vez a acciones no
idempotentes de larga duración y deja los flujos Express, al menos una vez, para
alto volumen ya idempotente \parencite{awsstepfunctions}. En esa última ruta la
plataforma no gestiona la idempotencia por el consumidor: uno que no lo sea
devuelve resultados incorrectos al primer reintento.
% TODO (Valentina)

%% ---------------------------------------------------------------------------
%% 2.5                                                    Responsable: Isidora
%% Portabilidad, aislamiento y arranque submilisegundo. WASI desde W3C /
%% Bytecode Alliance; Fermyon Spin / SpinKube, wasmCloud, Wasmtime.
%% Candidatos "otros" de edge y open source -> pasar a Claudio y Daniel (3.1).
%% ---------------------------------------------------------------------------
\subsection{WebAssembly y WASI como entorno de borde}

% TODO (Isidora)

%% ---------------------------------------------------------------------------
%% 2.6                                                    Responsable: Fabián
%% Tiempos máximos de ejecución, tamaño de carga útil, manejo de estado,
%% estrategias de portabilidad, lock-in. Se apoya en la Tabla 3.2 (límites)
%% y en las alternativas open source de 3.1 (referencias cruzadas).
%% ---------------------------------------------------------------------------
\subsection{Límites técnicos y dependencia del proveedor}

% TODO (Fabián)

\clearpage


%%%%%%%%%%%%%%%%%%%%%%%%%%%%%%%%%%%%%%%%%%%%%%%%%%%%%%%%%%%%%%%%%%%%%%%%%%%%%%%%
%% >>> 3. ANÁLISIS COMPARATIVO DE TECNOLOGÍAS Y PRODUCTOS
%% Entregable de la ficha: cuadro comparativo de al menos 6 plataformas con
%% sus límites técnicos. Aquí se levantan TODOS los precios del informe.
%% IA: prohibida en el análisis comparativo y la justificación de criterios
%% (punto 6.1).
%%%%%%%%%%%%%%%%%%%%%%%%%%%%%%%%%%%%%%%%%%%%%%%%%%%%%%%%%%%%%%%%%%%%%%%%%%%%%%%%

\section{Análisis comparativo de tecnologías y productos}

%% ---------------------------------------------------------------------------
%% 3.1                                          Responsables: Claudio y Daniel
%% Presupuesto: 0,75 página / ~300 palabras + tabla.
%% OBLIGATORIO (punto 4): la fila "Tecnologías, productos y proveedores a
%% comparar" de la ficha tiene CINCO listas que terminan en "y otros que el
%% grupo debe identificar". Por cada lista: al menos DOS alternativas que no
%% figuren, con qué aportan y por qué merecen entrar en la comparación. Si
%% para alguna lista no hay más, declararlo expresamente y documentar la
%% búsqueda (fuentes, criterios de descarte, fecha) -> detalle en Anexo D.
%% Las que califiquen entran como filas de la Tabla 3.2.
%% Listas: FaaS · contenedores sin servidor · cómputo en el borde · código
%% abierto · datos sin servidor.
%% ---------------------------------------------------------------------------
\subsection{Alternativas adicionales}
\label{sec:alternativas}

% TODO (Claudio y Daniel): dos párrafos + tabla.

\tabla{Alternativas adicionales incorporadas a la comparación}{alternativas-adicionales}{|p{3.2cm}|p{3.4cm}|p{7.2cm}|}{%
    \filaTabla{\textbf{Lista de la ficha} & \textbf{Alternativa} & \textbf{Qué aporta a la comparación}}
    \filaTabla{Funciones como servicio & & }
    \filaTabla{Funciones como servicio & & }
    \filaTabla{Contenedores sin servidor & & }
    \filaTabla{Contenedores sin servidor & & }
    \filaTabla{Cómputo en el borde & & }
    \filaTabla{Cómputo en el borde & & }
    \filaTabla{Código abierto & & }
    \filaTabla{Código abierto & & }
    \filaTabla{Datos sin servidor & & }
    \filaTabla{Datos sin servidor & & }
}

%% ---------------------------------------------------------------------------
%% 3.2                     Responsables: Valentina (criterios), Rodolfo (AWS y
%%                         Azure; consolida y escribe el análisis), Claudio y
%%                         Daniel (Google, Cloudflare, Fastly, Deno/Vercel/
%%                         Netlify, open source).
%% Presupuesto: 1,25 página / ~250 palabras + tabla (idealmente en una página
%% horizontal: ver \begin{landscape} si hace falta, paquete pdflscape).
%% Criterios: costo, latencia, complejidad operacional (las tres dimensiones
%% del alcance de la ficha), justificados en un párrafo + tabla chica.
%% Cuadro: mínimo 6 plataformas de más de una categoría + las alternativas
%% de 3.1 que califiquen. Columnas de límites y de precio.
%% PRECIOS (punto 5): producto, moneda, región, fecha de consulta, lista o
%% cotización. Verificados en la calculadora oficial. Sin fecha no vale.
%% "Solo por cotización" si el proveedor no publica precio.
%% ---------------------------------------------------------------------------
\subsection{Criterios y cuadro comparativo de plataformas}
\label{sec:cuadro}

Para evaluar las alternativas tecnológicas frente al modelo de referencia de instancias siempre encendidas, se establecen tres criterios estructurados que capturan las dimensiones técnica, económica y operativa requeridas por el caso de licitación:

\tabla{Criterios de evaluación aplicados al cuadro comparativo}{criterios}{|p{2.8cm}|p{5.4cm}|p{5.6cm}|}{%
    \filaTabla{\textbf{Criterio} & \textbf{Por qué importa} & \textbf{Cómo se mide en el cuadro}}
    \filaTabla{Costo y modelo de cobro & Determina la viabilidad económica en función del tráfico y el riesgo financiero de pisos mínimos o sobrecostos ocultos. & Unidad métrica de cómputo (vCPU-s, GB-s, tiempo de CPU), piso fijo mensual y precio de lista unitario fechado (\$) en región de referencia.}
    \filaTabla{Latencia y aislamiento & Condiciona el tiempo de respuesta frente a ráfagas y la frontera de seguridad multi-inquilino entre cargas de trabajo. & Tecnología de sandbox declarada (microVM, contenedor, isolate V8, Wasm), mitigaciones de arranque en frío y su impacto en facturación.}
    \filaTabla{Complejidad operacional & Cuantifica la carga administrativa y de mantenimiento transferida al equipo de desarrollo frente al proveedor. & Delimitación de responsabilidades sobre lista cerrada estandarizada, límites estrictos que fuerzan rediseño y necesidad de externalizar estado.}
}

El análisis comparativo se estructura en doce plataformas representativas que cubren los cuatro modelos de ejecución analizados, complementadas por el panorama de código abierto. Todas las tarifas y cuotas corresponden a documentación oficial consultada el 20 de septiembre de 2026.

\begin{landscape}
\subsubsection*{Tabla 3.2a: Criterios arquitectónicos, operativos y económicos}
\label{tab:cuadro-3-2a}

{\scriptsize
\begin{longtable}{|p{2.4cm}|p{1.8cm}|p{2.8cm}|p{3.2cm}|p{3.6cm}|p{1.2cm}|p{2.4cm}|p{3.0cm}|p{3.2cm}|}
\hline
\textbf{Plataforma} & \textbf{Modelo} & \textbf{Aislamiento} & \textbf{Unidad de cobro y piso} & \textbf{Precio de lista fechado} & \textbf{Cold start} & \textbf{Mitigación (factura)} & \textbf{Qué opera el equipo} & \textbf{Estado / Límite rediseño} \\ \hline
\endhead

AWS Lambda & FaaS & MicroVM Firecracker (cgroups/chroot por función) & Solicitudes (\$0.20/M) + GB-s reloj. Sin piso; escala a cero. & us-east-1, USD, 20-09-2026: \$0.0000166667/GB-s x86; Free 1M req/400k GB-s. Lista & (4.1) & Concurrencia aprovisionada: SÍ (\$0.0000041667/GB-s). SnapStart sin cargo & runtime, despliegue, observabilidad & SÍ obligatorio (S3/DynamoDB). Límite: 900~s máx, 6~MB payload sync \\ \hline

Azure Functions (Flex Consumption) & FaaS & Contenedor / microVM en entorno aislado por app & Solicitud (\$0.40/M) + GB-s (\$0.000026). Sin piso; escala a cero. & East US, USD, 20-09-2026: \$0.40/M req + \$0.000026/GB-s; Free 250k req / 100k GB-s. Lista & s/m & Instancias Always Ready: SÍ (\$0.000004/GB-s base + \$0.000016 ejec) & runtime, despliegue, observabilidad & SÍ obligatorio (Blob/Cosmos). Límite: HTTP máx 230~s (Load Balancer) \\ \hline

Google Cloud Run functions & FaaS & gVisor (gen 1) o microVM (gen 2) sobre Cloud Run & vCPU-s, GiB-s y req. Sin piso; escala a cero. Cloud Build/Artifact aparte. & us-east4 (Tier 1), USD, 20-09-2026: CPU \$0.000024/s, RAM \$0.0000025/s, req \$0.40/M. Free tier Tier 1. Lista & s/m & Min instances y Startup CPU boost: SÍ se facturan al arranque & runtime, despliegue, observabilidad & SÍ obligatorio. Límite: 9~min en eventos frente a 60~min en HTTP \\ \hline

Google Cloud Run (servicios) & Contenedor & Sandbox gVisor (gen 1) o microVM Linux (gen 2) & vCPU-s y GiB-s (100~ms) + req. Sin piso mensual; escala a cero. & us-east4 (Tier 1), USD, 20-09-2026: CPU \$0.000024/s, idle \$0.0000025/s, RAM \$0.0000025/s. Lista & (4.1) & Min instances (idle) y CPU boost: SÍ se facturan & runtime, red (VPC/SSL), despliegue, observabilidad & SÍ (disco en RAM). Límite: 60~min máx; tareas largas van a Jobs (168~h) \\ \hline

AWS Fargate (ECS) & Contenedor & MicroVM dedicada por tarea (Firecracker en AWS) & vCPU-h y GB-h por segundo (mín 1~min). Sin piso; escala a cero. & us-east-1, USD, 20-09-2026: \$0.04048/vCPU-h, \$0.004445/GB-h. Fargate Spot desc. 70\%. Lista & s/m & No aplica a nivel función: mantener tareas mínimas activas completas & runtime, red (VPC/ALB), despliegue, observabilidad & Efímero 20~GB; persistencia vía EFS/EBS. Sin timeout de plataforma \\ \hline

Azure Container Apps & Contenedor & Contenedor en clúster K8s/Envoy dedicado por entorno & vCPU-s y GiB-s activos. Sin piso; escala a cero. & East US, USD, 20-09-2026: vCPU \$0.000024/s, idle \$0.000003/s, RAM \$0.000003/s. Free 180k vCPU-s. Lista & s/m & Min replicas en reposo: SÍ se facturan a tarifa idle reducida & runtime, red (VNet), despliegue, observabilidad & SÍ (Azure Files). Límite: 300~s máx en HTTP; batch requiere Jobs \\ \hline

Cloudflare Workers (Paid) & Borde & Isolates V8 multi-tenant en proceso único (espacio virtual común) & Solicitudes (\$0.30/M) + tiempo CPU (\$0.02/M ms). \textbf{Piso: \$5.00/mes cuenta.} & global, USD, 20-09-2026: \$5/mes base (10M req, 30M CPU-ms); reloj sin cobro. Free 100k req/d. Lista & (4.1) & Estructural: isolate en proceso vivo ($<$5~ms); snapshot Wasm. Sin SKU aparte & runtime, despliegue, observabilidad & SÍ obligatorio (DO, KV, D1, R2). Límite: 128~MB RAM y 5~min CPU máx \\ \hline

Fastly Compute & Borde & WebAssembly (WASI) sobre Wasmtime; sandbox por req. & Requests (\$0.50--0.20/M) + vCPU-ms (\$0.05--0.02/M ms). Sin piso. & global, USD, 20-09-2026: tramos uso actual. Free 10M req, 100M ms. Tráfico entrega aparte. Lista & s/m & Estructural: instanciación Wasm en microsegundos. Sin SKU aparte & runtime, despliegue, observabilidad & SÍ (Config/KV/Secret Store). Límite: 50~ms CPU por req., 128~MB heap \\ \hline

AWS Lambda@Edge / CloudFront & Borde & L@E: microVM en PoP edge; CF Functions: proceso liviano JS & L@E: \$0.60/M + \$0.00005001/GB-s. CF Func: \$0.10/M (2M free). Sin piso. & global, USD, 20-09-2026: precios de lista publicados en CDN CloudFront & s/m & CF Func submilisegundo estructural. L@E sin concurrencia aprovisionada & runtime, despliegue (asociación a CDN) & SÍ (sin disco local). Límite: CF Func 1~ms CPU; L@E 5~s viewer / 30~s origin \\ \hline

EC2 reservada (ECS t3.medium) & Siempre encendido & VM dedicada Nitro/KVM; sin aislamiento por invocación & Instancia-hora. \textbf{Piso: mes completo (\$30.37/mes On-Demand). No escala a cero.} & us-east-1, USD, 20-09-2026: On-Demand \$0.0416/h; Reserva 1a \$0.026/h (\$18.98/m); 3a \$0.017/h (\$12.41/m). & n/a & No aplica: proceso residente en memoria absorbe tráfico de inmediato & SO/host, Docker, capacidad, escalado, red, despliegue, observabilidad & Estado en memoria/disco host (EBS). Sin timeout de plataforma \\ \hline

Oracle OCI Functions (3.1) & FaaS & Contenedor gestionado / microVM OCI (motor Fn Project) & Invocaciones (\$0.20/M) + GB-s (\$0.00001417/s). Concurrency idle 25\%. Sin piso. & US East, USD, 20-09-2026: Always Free 2M req + 400k GB-s/mes permanentes. Lista & s/m & Concurrencia aprovisionada: SÍ (25\% tarifa en idle; sin recargo activo) & runtime, despliegue, observabilidad & SÍ (Object Storage, DB). Límite: 300~s sync; modo detached hasta 3600~s \\ \hline

Azion Edge Functions (3.1) & Borde & Isolates V8 sobre Edge Runtime distribuido propio & Compute (\$0.072/1000~s CPU) + req (\$0.30/M). Sin piso en Starter. & Brasil/LATAM/global, USD, 20-09-2026: Free 8~h compute + 3M req/mes. Lista & s/m & Aislamiento liviano V8 ($<$2~s cold start doc, típico ms). Sin SKU & runtime, despliegue, observabilidad & SÍ (Edge Storage/SQL). Límite: \textbf{2~s CPU por req} (wall 5~min); isolate 512~MB \\ \hline
\end{longtable}
}

\clearpage

\subsubsection*{Tabla 3.2b: Límites técnicos y tarifas de egreso de datos}
\label{tab:cuadro-3-2b}

{\scriptsize
\begin{longtable}{|p{2.6cm}|p{3.2cm}|p{3.4cm}|p{3.2cm}|p{3.6cm}|p{4.2cm}|}
\hline
\textbf{Plataforma} & \textbf{Tiempo máx. ejecución} & \textbf{Memoria / vCPU} & \textbf{Payload in / out} & \textbf{Concurrencia / escalado} & \textbf{Egreso de datos (precio fechado)} \\ \hline
\endhead

AWS Lambda & 900~s (15~min) & 128 a 10.240~MB (1~vCPU a 1.769~MB, máx 6~vCPU) & Sync: 6~MB. Async: 256~KB & 1.000 concurrentes regional default (ampliable). Escala a cero & 100~GB/mes gratis; luego \$0.09/GB (primeros 10~TB us-east-1). 20-09-2026 \\ \hline

Azure Functions (Flex) & HTTP: 230~s (techo Load Balancer); otros triggers 30~min (máx unbounded) & 512~MB (0.5~vCPU), 2048~MB (1~vCPU), 4096~MB (2~vCPU) & HTTP 100~MB en gateway; Storage Queues 4~MB & Hasta 1.000 instancias por app. Escala a cero & 100~GB/mes gratis; luego \$0.087/GB (East US). 20-09-2026 \\ \hline

Google Cloud Run functions & HTTP hasta 60~min; eventos hasta 9~min (API v2) & Default 256~MiB; máx 16~GiB / 4~vCPU según tabla funciones & HTTP/1 32~MiB req/resp. Sin límite bajo HTTP/2 & Hasta 1.000 req concurrentes por instancia. Escala a cero & Premium Tier: 1~GiB gratis; luego \$0.12/GB hacia Norteamérica. 20-09-2026 \\ \hline

Google Cloud Run (servicios) & 300~s default; máx 3600~s (60~min). Jobs 168~h & Default 512~MiB (mín 128~MiB, máx 32~GiB); vCPU default 1 (0.08 a 8) & HTTP/1 32~MiB. Sin límite con streaming HTTP/2 & Instancia: default 80, máx 1.000. Revisión: máx 100 inst. default. Escala a cero & Premium Tier: 1~GiB gratis; luego \$0.12/GB (1--1.024~GiB). Interno sin cargo. 20-09-2026 \\ \hline

AWS Fargate (ECS) & Sin límite de plataforma (proceso continuo) & 0.25~vCPU (0.5--2~GB) a 32~vCPU (hasta 244~GB en doc ECS) & Sin límite de plataforma (ALB máx 100~MB cuerpo) & Cuota default 1.000 tareas regional. Escala a cero si desired=0 & Salida vía NAT Gateway (\$0.045/GB + \$0.045/h) + egreso EC2 (\$0.09/GB). 20-09-2026 \\ \hline

Azure Container Apps & HTTP 300~s máx; Jobs sin límite (hasta 7 días) & 0.25~vCPU / 0.5~GiB a 4.0~vCPU / 8.0~GiB en Consumption & HTTP hasta 100~MB soportados vía ingress Envoy & Hasta 300 réplicas por app (default 10--30). Escala a cero & 100~GB/mes gratis; luego \$0.087/GB hacia internet. 20-09-2026 \\ \hline

Cloudflare Workers (Paid) & CPU: 5~min máx (\textbf{30~s default}). Reloj: \textbf{SIN LÍMITE} en HTTP & 128~MB por isolate, NO configurable. vCPU no expuesta & URL 16~KB; cabeceras 128~KB; body req según plan (100~MB a 5~GB) & Sin límite numérico publicado; 10.000 subsolicitudes en Paid. Escala a cero & \textbf{SIN CARGO declarado:} egreso y ancho de banda incluidos en Workers Paid. 20-09-2026 \\ \hline

Fastly Compute & CPU: \textbf{50~ms por ejecución}. Reloj: 2~min (60~s trial) & 1~MB stack y 128~MB heap por ejecución, no configurable & Cabeceras 128~KB; URL 8~KB; caché 100~MB; paquete Wasm 100~MB & Instancia aislada por req. Encadenamiento máx 20 saltos y 6 servicios & \textbf{Entrega (Full Site Delivery) aparte:} 100~GB free; luego \$0.12/GB (US/EU), \$0.19 (LATAM). 20-09-2026 \\ \hline

AWS Lambda@Edge / CF Func & CF Func: \textbf{1~ms CPU}. L@E: 5~s viewer / 30~s origin & CF Func: 2~MB fijos. L@E: 128~MB (viewer) o 128--10.240~MB (origin) & CF Func: sin cuerpo. L@E: 40~KB viewer, 1~MB origin (40~MB stream) & Integrado en escala CDN global CloudFront. Escala a cero & Tarifa CloudFront: 1~TB/mes gratis; luego \$0.085/GB primer tramo en Norteamérica. 20-09-2026 \\ \hline

EC2 reservada (ECS t3.medium) & Sin límite de plataforma (ejecución continua) & Fijo de instancia: 2~vCPU y 4.0~GiB de RAM dedicados & Sin límite de plataforma; gobernado por memoria de host y proxy & Capacidad física de t3.medium; autoescalado por ASG (no a cero) & Salida estándar EC2: 100~GB/mes gratis; luego \$0.09/GB (us-east-1). 20-09-2026 \\ \hline

Oracle OCI Functions (3.1) & Sync: 300~s (\textbf{default 30~s}). Detached: 5 a \textbf{3.600~s (60~min)} & Opciones fijas: 128, 256, 512, 1024, 2048 y 3072~MB & Solicitud y respuesta HTTP: 6~MB cada una & 60 concurrentes por AD default (ampliable). Escala a cero & \textbf{Primeros 10~TB/mes gratis hacia internet}; luego \$0.0085/GB. 20-09-2026 \\ \hline

Azion Edge Functions (3.1) & CPU: \textbf{2~s máx por req}. Reloj wall-clock: máx \textbf{5~min (300~s)} & Hasta 512~MB por isolate de función. vCPU no expuesta & Args JSON 100~KB. Código: UI 6~MB, API 20~MB. Sub-requests: 50 & Autoescalado sobre red distribuida de PoPs edge. Escala a cero & Red Edge por GB (Brasil/LATAM \$0.06/GB en inicio; 50~GB gratis). 20-09-2026 \\ \hline
\end{longtable}
}
\end{landscape}

\subsubsection*{Tabla 3.2c: Ecosistema de plataformas de código abierto}
\label{tab:cuadro-3-2c}

{\scriptsize
\begin{longtable}{|p{2.6cm}|p{2.8cm}|p{2.6cm}|p{1.6cm}|p{4.2cm}|p{3.2cm}|}
\hline
\textbf{Proyecto} & \textbf{Modelo} & \textbf{Aislamiento} & \textbf{Escala 0 default} & \textbf{Defaults configurables (timeout / concurrencia)} & \textbf{Acreditación CNCF u otra} \\ \hline
\endhead

Knative Serving & Contenedores K8s & Contenedor (pod K8s) & \textbf{Sí (30~s gracia)} & Timeout 300~s default (máx 600~s). Concurrencia 0 (ilimitada default, máx 1000) & \textbf{Graduado CNCF (11-09-2025)}. Categoría: Scheduling \& Orchestration \\ \hline

OpenFaaS & FaaS sobre K8s / host & Contenedor (gVisor opcional en Pro) & No (opt-in) & Timeouts en gateway y función configurables. JetStream ack\_wait 60~s. Soporta \texttt{faasd} sin K8s & \textbf{No es CNCF}. OpenFaaS Ltd. Licencias CE (EULA) y comercial Pro/Edge \\ \hline

KEDA & Autoescalador K8s & No aplica (escala otros) & \textbf{Sí} & No aplica en ejecución HTTP. Sondeo default 30~s (\texttt{pollingInterval}) & \textbf{Graduado CNCF (22-08-2023)}. Scheduling \& Orchestration \\ \hline

Fermyon SpinKube & Serverless Wasm en K8s & WebAssembly (\texttt{containerd-shim-spin}) & No nativo (delega KEDA) & Timeouts gobernados por pod de K8s y triggers de Spin. Wasm en nodo sin contenedor & \textbf{Sandbox CNCF (21-01-2025)}. Runtime Spin en org \texttt{spinframework} \\ \hline

wasmCloud & Componentes Wasm & WebAssembly (Wasmtime Component Model) & \textbf{Sí (por diseño)} & Concurrencia ilimitada default (\texttt{wash scale}). Control timeout 2000~ms default & \textbf{Incubating CNCF (08-11-2024)}. Scheduling \& Orchestration \\ \hline

Fission (3.1) & FaaS sobre K8s & Contenedor (pods precalentados o dedicados) & Configurable por función & Router timeout 300~s default. Warm pools (\texttt{poolmgr}, 3 pods) vs escala a cero (\texttt{newdeploy}) & \textbf{No es CNCF}. Proyecto comunitario activo (v1.27.0, 2026) \\ \hline
\end{longtable}
}

\subsubsection*{Notas metodológicas y de trazabilidad al pie de las tablas}
\begin{itemize}
    \item \textbf{Nota (a) --- Complejidad operacional («Qué opera el equipo»):} La lista cerrada corresponde a la matriz estandarizada construida por el grupo a partir de las matrices de responsabilidad de AWS, Azure y Google Cloud, y el Whitepaper de Serverless del CNCF. Los proveedores de borde no publican una matriz formal; las responsabilidades son una inferencia técnica documentada desde sus contratos de arquitectura.
    \item \textbf{Nota (b) --- Arranque en frío:} La marca \texttt{(4.1)} remite a los ensayos de laboratorio del grupo (Vicente Arratia). Para las demás plataformas se declara \texttt{sin medición propia} (\texttt{s/m}) para no distorsionar la comparación con benchmarks heterogéneos de proveedores o literatura. La instancia siempre encendida declara \texttt{no aplica} (\texttt{n/a}).
    \item \textbf{Nota (c) --- Lambda@Edge y CloudFront Functions:} Se consolidan en fila conjunta por representar la oferta escalonada de AWS sobre CDN, diferenciando manipulación rápida de cabeceras (1~ms CPU) de cómputo regional (hasta 30~s).
    \item \textbf{Nota (d) --- Respaldo teórico del aislamiento en el borde:} Cloudflare Workers y Fastly Compute utilizan aislamiento a nivel de lenguaje (\emph{V8 Isolates} y \emph{WebAssembly/WASI}) en un mismo espacio de memoria virtual para ofrecer arranque sub-milisegundo. Como demuestran Schwarzl et al.\ \parencite{schwarzl2022spectre} en \emph{Robust and Scalable Process Isolation against Spectre in the Cloud}, este enfoque optimiza la densidad a costa de exposición a canales laterales microarquitectónicos (\emph{Spectre-PHT}), lo que obliga a incorporar defensas probabilísticas activas (\emph{DyPrIs}).
\end{itemize}

%% ---------------------------------------------------------------------------
%% Análisis comparativo de plataformas (Sección 3.2)
%% ---------------------------------------------------------------------------
El análisis transversal del cuadro comparativo revela que la promesa teórica de «escala a cero y pago estricto por uso» experimenta tensiones significativas cuando se confronta con las exigencias operativas reales. En primer lugar, la mitigación de arranques en frío en FaaS de propósito general introduce costos fijos encubiertos: tanto \emph{Provisioned Concurrency} en AWS Lambda (\$0,0000041667/GB-s reservado más \$0,0000097222/GB-s activo) como las instancias \emph{Always Ready} en Azure Functions Flex (\$0,000004/GB-s de reserva continua) transforman la facturación variable en un piso mensual análogo al de una instancia aprovisionada. En cargas con tráfico sostenido o ráfagas regulares (como la recepción continua de telemetría de grúas y contenedores en TERABYTE), este sobrecosto aproxima el gasto mensual al punto de equilibrio de una instancia dedicada reservada (EC2 \texttt{t3.medium} a \$18,98/mes, Fila~10), donde el proceso residente en memoria elimina el retardo por diseño. Asimismo, el cómputo en el borde evidencia una bifurcación económica y técnica crítica: ejecutar lógica en la periferia mediante microVMs regionales (AWS Lambda@Edge) impone un costo de cómputo de \$0,00005001/GB-s, exactamente 3,0 veces superior al de Lambda estándar (\$0,0000166667/GB-s), además de carecer de capa gratuita y trasladar tarifas de distribución CDN sin posibilidad de concurrencia aprovisionada. En contraste, las soluciones basadas en aislamiento a nivel de lenguaje (Cloudflare Workers y Fastly Compute) eliminan el arranque en frío mediante procesos e isolates compartidos sub-milisegundo, pero exigen una renuncia arquitectónica severa: operan bajo estrictos límites de CPU (50~ms por solicitud en Fastly y 2~s en Azion), carecen de sistemas de archivos convencionales y trasladan la persistencia a bases externas mediante llamadas de red adicionales.

En segundo lugar, la complejidad operacional no desaparece en las arquitecturas serverless, sino que se redistribuye entre la infraestructura y el diseño de la aplicación. La transición desde FaaS hacia contenedores sin servidor (AWS Fargate y Azure Container Apps) amplía la flexibilidad de empaquetado y elimina las cuotas de tiempo (soportando procesos batch continuos o tareas de hasta 7 días), pero transfiere al equipo responsabilidades de ingeniería de red considerables: requiere provisionar subredes VPC privadas, balanceadores de carga de aplicación (ALB) y pasarelas NAT (cuyo costo base de \$0,045/hora más \$0,045/GB suele exceder el cómputo en volúmenes medios), además de sintonizar manualmente políticas de autoescalado (KEDA o métricas de CPU). Esta fricción se acentúa por límites duros no configurables de la infraestructura subyacente que contradicen las cuotas promocionadas por los proveedores: en Azure Functions, si bien la plataforma publicita tiempos de ejecución de hasta 60~minutos o modo no acotado (\emph{unbounded}), las funciones disparadas por HTTP enfrentan un techo infranqueable de 230~segundos impuesto por el \emph{idle timeout} del Azure Load Balancer, obligando a rediseñar procesos analíticos hacia patrones asíncronos desacoplados mediante colas o \emph{Durable Functions}. Finalmente, el ecosistema serverless introduce un riesgo latente de dependencia y obsolescencia que valida la necesidad de evaluar alternativas portables (Sección~\ref{sec:alternativas}): el anuncio oficial de AWS cerrando App Runner a nuevos clientes en 2026 demuestra que las abstracciones gestionadas de alta conveniencia están sujetas a discontinuación unilateral del proveedor, consolidando a las plataformas basadas en contenedores abiertos y estándares de la CNCF (Knative, Fn Project, Fission) como salvaguardas indispensables para la soberanía arquitectónica institucional.

\clearpage


%%%%%%%%%%%%%%%%%%%%%%%%%%%%%%%%%%%%%%%%%%%%%%%%%%%%%%%%%%%%%%%%%%%%%%%%%%%%%%%%
%% >>> 4. CASO DE ESTUDIO E IMPLEMENTACIÓN PRÁCTICA
%% Dos entregables de la ficha (PoC y modelo de costos) más la aplicación al
%% caso de licitación. Los tres comparten escenario: la PoC corre en
%% plataformas que están en la Tabla 3.2; el modelo cotiza una de esas
%% plataformas con los precios de la Tabla 3.2; el caso TERABYTE usa el
%% punto de equilibrio de 4.2 y la latencia de 4.1.
%% Esto es el "caso práctico" del 70 % de la nota.
%%%%%%%%%%%%%%%%%%%%%%%%%%%%%%%%%%%%%%%%%%%%%%%%%%%%%%%%%%%%%%%%%%%%%%%%%%%%%%%%

\section{Caso de estudio e implementación práctica}

%% ---------------------------------------------------------------------------
%% 4.1                                                    Responsable: Vicente
%% Presupuesto: 1,25 página / ~350 palabras + 1 gráfico + 1 tabla resumen.
%% Ficha: "prueba de concepto documentada (código y mediciones) de arranque
%% en frío y latencia, O BIEN análisis de mediciones publicadas con su
%% metodología". Se hace la propia: función trivial en 2 o 3 plataformas de
%% la Tabla 3.2 (Lambda, Cloud Run, Cloudflare Workers; free tier), script
%% que mide cold start vs. warm, N invocaciones, misma región si es posible.
%% En el cuerpo: metodología (1 párrafo), gráfico, tabla resumen (p50/p95
%% cold y warm por plataforma), interpretación y limitaciones.
%% Código, paso a paso y datos crudos -> Anexo B. Código auxiliar generado
%% con IA se identifica como tal y se declara.
%% ---------------------------------------------------------------------------
\subsection{Prueba de concepto: arranque en frío y latencia}
\label{sec:poc}

% TODO (Vicente): metodología.

% \imagen{assets/images/poc-coldstart.png}{Latencia de arranque en frío y en caliente por plataforma}{fig:poc}

\tabla{Resumen de mediciones de la prueba de concepto}{poc-resumen}{|p{3cm}|p{2.2cm}|p{2.2cm}|p{2.2cm}|p{2.2cm}|p{2.4cm}|}{%
    \filaTabla{\textbf{Plataforma} & \textbf{Cold p50 (ms)} & \textbf{Cold p95 (ms)} & \textbf{Warm p50 (ms)} & \textbf{Warm p95 (ms)} & \textbf{N / región / fecha}}
    \filaTabla{ & & & & & }
    \filaTabla{ & & & & & }
    \filaTabla{ & & & & & }
}

% TODO (Vicente): interpretación y limitaciones.

%% ---------------------------------------------------------------------------
%% 4.2                                                    Responsable: Francisca
%% Presupuesto: 1 página / ~300 palabras + tabla de supuestos + 1 gráfico.
%% Ficha: "modelo de costos que compare una arquitectura sin servidor con una
%% de contenedores siempre encendidos, identificando el volumen de
%% equilibrio en solicitudes por mes".
%% En el cuerpo: supuestos (duración media, memoria, egress, tamaño de
%% respuesta, escenarios de carga), gráfico costo mensual vs. solicitudes/mes
%% con las dos curvas y el cruce, el número de equilibrio.
%% Planilla completa -> Anexo C. Precios: solo los de la Tabla 3.2.
%% ---------------------------------------------------------------------------
\subsection{Modelo de costos y volumen de equilibrio}
\label{sec:costos}

% TODO (Francisca)

\tabla{Supuestos del modelo de costos}{supuestos-costos}{|p{5cm}|p{3.5cm}|p{5.5cm}|}{%
    \filaTabla{\textbf{Variable} & \textbf{Valor} & \textbf{Fuente / justificación}}
    \filaTabla{Duración media por solicitud (ms) & & }
    \filaTabla{Memoria asignada (MB) & & }
    \filaTabla{Tamaño medio de respuesta (KB) & & }
    \filaTabla{Plataforma serverless cotizada & & Tabla \ref{tab:cuadro-comparativo}}
    \filaTabla{Configuración always-on cotizada & & }
    \filaTabla{Región y fecha de precios & & }
}

% \imagen{assets/images/punto-equilibrio.png}{Costo mensual en función del volumen de solicitudes: serverless frente a contenedores siempre encendidos}{fig:equilibrio}

%% ---------------------------------------------------------------------------
%% 4.3                                                    Responsable: Fabián
%% Presupuesto: 0,75 página / ~350 palabras.
%% "Aplicación concreta al caso de licitación del propio grupo" (punto 4, lo
%% que se valora) + "Vínculo con la propuesta técnico-económica" de la ficha:
%% paso de CAPEX a OPEX puro y el riesgo de una factura que crece con el
%% éxito del sistema. Usar el volumen estimado del proyecto TERABYTE, el
%% punto de equilibrio de 4.2 y la latencia de 4.1 para decir qué modelo
%% conviene y cómo se refleja en la propuesta. No redescribir el proyecto.
%% ---------------------------------------------------------------------------
\subsection{Aplicación al caso de licitación de TERABYTE}
\label{sec:caso}

% TODO (Fabián)

\clearpage


%%%%%%%%%%%%%%%%%%%%%%%%%%%%%%%%%%%%%%%%%%%%%%%%%%%%%%%%%%%%%%%%%%%%%%%%%%%%%%%%
%% >>> 5. TENDENCIAS Y RIESGOS                            Responsable: Fabián
%% Presupuesto: 1 página / ~450 palabras.
%% Sección exigida por las Pautas del curso ("Tendencias y evolución
%% futura"), escrita con el criterio de las Indicaciones ("evidencia
%% contradictoria bien tratada"). Cada tendencia va con su contraevidencia.
%% Máximo dos casos públicos de equipos que abandonaron serverless por costo
%% o complejidad, verificados en la fuente original (blogs de ingeniería:
%% admisibles declarándolo). Ninguna frase sin fuente.
%% IA: prohibida en la discusión crítica de la evidencia (punto 6.1).
%%%%%%%%%%%%%%%%%%%%%%%%%%%%%%%%%%%%%%%%%%%%%%%%%%%%%%%%%%%%%%%%%%%%%%%%%%%%%%%%

\section{Tendencias y riesgos}

% TODO (Fabián)

\clearpage


%%%%%%%%%%%%%%%%%%%%%%%%%%%%%%%%%%%%%%%%%%%%%%%%%%%%%%%%%%%%%%%%%%%%%%%%%%%%%%%%
%% >>> CONCLUSIONES Y RECOMENDACIÓN                       Responsable: todo el grupo
%% Presupuesto: 1 página / ~450 palabras.
%% Responder literalmente la pregunta de la ficha: en qué condiciones de
%% carga conviene cada modelo (FaaS, contenedor serverless, edge, always-on)
%% en costo, latencia y complejidad operacional. Terminar con UNA
%% recomendación con la que el grupo se compromete (punto 4).
%% No resumir el informe.
%% IA: prohibida (punto 6.1). Todos deben poder defenderla oralmente.
%%%%%%%%%%%%%%%%%%%%%%%%%%%%%%%%%%%%%%%%%%%%%%%%%%%%%%%%%%%%%%%%%%%%%%%%%%%%%%%%

\phantomsection
\addcontentsline{toc}{section}{Conclusiones y recomendación}
\section*{Conclusiones y recomendación}

% TODO (grupo)

\clearpage


%%%%%%%%%%%%%%%%%%%%%%%%%%%%%%%%%%%%%%%%%%%%%%%%%%%%%%%%%%%%%%%%%%%%%%%%%%%%%%%%
%% >>> BIBLIOGRAFÍA                                       Responsable: todos
%% Entradas en referencias.bib; citar con \parencite{clave}.
%% Mínimo 10 fuentes académicas (papers revisados por pares, libros,
%% estándares) + fuentes técnicas (docs oficiales). Cada uno aporta 1 o 2
%% académicas de su sección.
%% Preferentes: docs oficiales, NIST/ISO/IETF/CNCF/OWASP/ASF, papers,
%% bibliotecas PUCV. Admisibles declarándolo: informes de industria (muestra
%% y metodología), blogs de ingeniería. NO admisibles como primarias:
%% comparativas de un proveedor sobre su categoría, sin autor, agregadores.
%% Toda referencia se abre y verifica en el original. Referencia inexistente
%% = error grave. Entradas con precios: incluir región y fecha de consulta
%% (campo note o urldate).
%%%%%%%%%%%%%%%%%%%%%%%%%%%%%%%%%%%%%%%%%%%%%%%%%%%%%%%%%%%%%%%%%%%%%%%%%%%%%%%%

\phantomsection
\addcontentsline{toc}{section}{Bibliografía}
\printbibliography[title={Bibliografía}]

\clearpage


%%%%%%%%%%%%%%%%%%%%%%%%%%%%%%%%%%%%%%%%%%%%%%%%%%%%%%%%%%%%%%%%%%%%%%%%%%%%%%%%
%% >>> ANEXO A. DECLARACIÓN DE USO DE INTELIGENCIA ARTIFICIAL
%%                                          Responsable: Daniel recopila;
%%                                          cada integrante declara lo suyo.
%% OBLIGATORIO aunque todos declaren nivel 0 (punto 6.3). Sin este anexo el
%% informe no se recibe conforme. Prevalece sobre la lámina de Pautas que
%% decía "Anexos (opcional)".
%% Numeración A1, A2, ... (macro \anexo, definida en el preámbulo; cada anexo
%% reinicia páginas, tablas y figuras con su propia letra).
%%%%%%%%%%%%%%%%%%%%%%%%%%%%%%%%%%%%%%%%%%%%%%%%%%%%%%%%%%%%%%%%%%%%%%%%%%%%%%%%

\anexo{A}{Anexo A: Declaración de uso de inteligencia artificial}

\textoPortadaFont{\textbf{Empresa:} N.\textordmasculine\ 6 --- TERABYTE\\}
\textoPortadaFont{\textbf{Tema:} TI-05 --- Serverless y computación en el borde (edge computing)\\}
\textoPortadaFont{\textbf{Fecha:} 21 de septiembre de 2026\\}

\bigskip

\subsection*{A. Declaración por sección}

%% Niveles (punto 6.2): 0 sin IA · 1 corrección y estilo · 2 colaboración e
%% ideación (esquemas, resumen de bibliografía, código de apoyo) · 3
%% generación sustancial. Niveles 2 y 3 exigen prompts (B) y evidencia (C).
%% Secciones donde NO se admite IA (punto 6.1): análisis comparativo y
%% criterios (3.2), conclusiones, párrafo de aporte, discusión crítica (5),
%% cuestionario, y cualquier cifra, cita o referencia.
%% La tabla incluye el cuestionario y la presentación porque el formulario
%% oficial los lista como filas.

{\scriptsize
\begin{longtable}{|p{5.6cm}|p{1.6cm}|p{3.4cm}|p{4.3cm}|}
    \hline
    \textbf{Sección del informe} & \textbf{Nivel (0-3)} & \textbf{Herramienta y versión} & \textbf{Evidencia adjunta (niveles 2 y 3)} \\
    \hline
    \endhead
    Resumen ejecutivo & & & \\ \hline
    Introducción y contexto (incl. párrafo de aporte) & & & \\ \hline
    1. Marco teórico y conceptual & & & \\ \hline
    2.1 Modelos de ejecución y aislamiento & & & \\ \hline
    2.2 Arranque en frío & & & \\ \hline
    2.3 Facturación por consumo y punto de equilibrio & & & \\ \hline
    2.4 Diseño orientado a eventos & & & \\ \hline
    2.5 WebAssembly y WASI & & & \\ \hline
    2.6 Límites técnicos y dependencia del proveedor & & & \\ \hline
    3.1 Alternativas adicionales & & & \\ \hline
    3.2 Criterios y cuadro comparativo & & & \\ \hline
    4.1 Prueba de concepto (incl. código auxiliar) & & & \\ \hline
    4.2 Modelo de costos & & & \\ \hline
    4.3 Aplicación al caso TERABYTE & & & \\ \hline
    5. Tendencias y riesgos & & & \\ \hline
    Conclusiones y recomendación & & & \\ \hline
    Bibliografía & & & \\ \hline
    Cuestionario de 30 preguntas & & & \\ \hline
    Presentación (diapositivas) & & & \\ \hline
\end{longtable}
}

\subsection*{B. Prompts utilizados (niveles 2 y 3)}

% TODO: transcribir los prompts efectivamente empleados, indicando sección.
% Si son muchos, hoja adjunta identificada o enlace.

\subsection*{C. Evidencia trazable adjunta (niveles 2 y 3)}

% TODO: enlaces a conversaciones exportadas + historial de versiones
% (Overleaf History / Google Docs / repositorio con commits). Verificar que
% el profesor pueda acceder.

\subsection*{D. Firma de los integrantes}

%% Texto del formulario oficial: "Declaro que la información consignada en
%% este formulario es completa y veraz, que las secciones señaladas en el
%% punto 6.1 de las indicaciones son de autoría humana y que puedo explicar
%% oralmente cualquier parte del trabajo entregado."
%% Cada integrante completa su fila. Nadie firma por otro.

Declaro que la información consignada en este formulario es completa y veraz, que las secciones señaladas en el punto 6.1 de las indicaciones son de autoría humana y que puedo explicar oralmente cualquier parte del trabajo entregado.

{\scriptsize
\begin{longtable}{|p{0.8cm}|p{6.2cm}|p{3.8cm}|p{4.1cm}|}
    \hline
    \textbf{N.\textordmasculine} & \textbf{Nombre completo} & \textbf{Nivel máximo declarado} & \textbf{Firma} \\
    \hline
    \endhead
    1 & Francisca Antonia Abarca Pezo & & \\ \hline
    2 & Vicente André Arratia Caroca & & \\ \hline
    3 & Isidora Valentina Cisternas Cisternas & & \\ \hline
    4 & Rodolfo Antonio Fernández Vera & & \\ \hline
    5 & Valentina Ignacia Guzmán Elgueta & & \\ \hline
    6 & Daniel Ignacio Miranda Vázquez & & \\ \hline
    7 & Matias Ignacio Reyes Pizarro & & \\ \hline
    8 & Fabián Andres Solís Piñones & & \\ \hline
    9 & Claudio Patricio Toledo Mac-lean & & \\ \hline
\end{longtable}
}

\clearpage


%%%%%%%%%%%%%%%%%%%%%%%%%%%%%%%%%%%%%%%%%%%%%%%%%%%%%%%%%%%%%%%%%%%%%%%%%%%%%%%%
%% >>> ANEXO B. CÓDIGO Y MEDICIONES DE LA PRUEBA DE CONCEPTO   Responsable: Vicente
%% Código de la función y del script de medición (entorno minted, ya cargado
%% por la clase; Overleaf debe compilar con -shell-escape), paso a paso de
%% despliegue, tabla de datos crudos o enlace al repositorio.
%%%%%%%%%%%%%%%%%%%%%%%%%%%%%%%%%%%%%%%%%%%%%%%%%%%%%%%%%%%%%%%%%%%%%%%%%%%%%%%%

\anexo{B}{Anexo B: Código y mediciones de la prueba de concepto}

% TODO (Vicente)
% \begin{minted}{python}
% ...
% \end{minted}

\clearpage


%%%%%%%%%%%%%%%%%%%%%%%%%%%%%%%%%%%%%%%%%%%%%%%%%%%%%%%%%%%%%%%%%%%%%%%%%%%%%%%%
%% >>> ANEXO C. PLANILLA DEL MODELO DE COSTOS               Responsable: Francisca
%% Tabla completa por escenario de carga (por ejemplo 10 mil, 100 mil, 1
%% millón y 10 millones de solicitudes/mes), fórmulas, y la cita de cada
%% precio con producto, moneda, región, fecha y tipo. Puede ir como tabla
%% LaTeX o como PDF externo (\pdfExterno{assets/pdf/...}{...}{...}).
%%%%%%%%%%%%%%%%%%%%%%%%%%%%%%%%%%%%%%%%%%%%%%%%%%%%%%%%%%%%%%%%%%%%%%%%%%%%%%%%

\anexo{C}{Anexo C: Planilla del modelo de costos}

% TODO (Francisca)

\clearpage


%%%%%%%%%%%%%%%%%%%%%%%%%%%%%%%%%%%%%%%%%%%%%%%%%%%%%%%%%%%%%%%%%%%%%%%%%%%%%%%%
%% >>> ANEXO D. BÚSQUEDA DOCUMENTADA DE ALTERNATIVAS   Responsables: Claudio y Daniel
%% Punto 4, "Declaración cuando no hay más": por cada una de las cinco
%% listas, fuentes consultadas, criterios de inclusión y descarte, fecha, y
%% candidatos descartados con su motivo (por ejemplo, producto descontinuado).
%%%%%%%%%%%%%%%%%%%%%%%%%%%%%%%%%%%%%%%%%%%%%%%%%%%%%%%%%%%%%%%%%%%%%%%%%%%%%%%%

\anexo{D}{Anexo D: Búsqueda documentada de alternativas}

% TODO (Claudio y Daniel)

\tabla{Búsqueda de alternativas por lista de la ficha}{busqueda-alternativas}{|p{2.8cm}|p{3.6cm}|p{3.2cm}|p{2.6cm}|p{1.8cm}|}{%
    \filaTabla{\textbf{Lista} & \textbf{Fuentes consultadas} & \textbf{Criterios de inclusión / descarte} & \textbf{Descartadas (motivo)} & \textbf{Fecha}}
    \filaTabla{Funciones como servicio & & & & }
    \filaTabla{Contenedores sin servidor & & & & }
    \filaTabla{Cómputo en el borde & & & & }
    \filaTabla{Código abierto & & & & }
    \filaTabla{Datos sin servidor & & & & }
}

\clearpage


%%%%%%%%%%%%%%%%%%%%%%%%%%%%%%%%%%%%%%%%%%%%%%%%%%%%%%%%%%%%%%%%%%%%%%%%%%%%%%%%
%% >>> DOCUMENTO ADJUNTO: CUESTIONARIO DE 30 PREGUNTAS
%%                              Responsable: todos (3-4 c/u); Daniel arma el
%%                              índice temático y ajusta la proporción.
%% Punto 3 de las Indicaciones:
%%   - 30 preguntas; formatos mezclados: selección múltiple, V/F, completar,
%%     respuesta corta;
%%   - cada una con respuesta correcta y breve justificación;
%%   - teoría y práctica; comprensión, aplicación y análisis crítico;
%%   - dificultad: 12 básicas, 12 intermedias, 6 avanzadas (40/40/20);
%%   - índice temático por sección del informe.
%% Redacción 100 % humana (punto 6.1). Numeración Q1, Q2, ... (la letra B
%% ya la usa el Anexo B).
%%%%%%%%%%%%%%%%%%%%%%%%%%%%%%%%%%%%%%%%%%%%%%%%%%%%%%%%%%%%%%%%%%%%%%%%%%%%%%%%

\anexo{Q}{Documento adjunto: Cuestionario de evaluación de conocimientos}

\subsection*{Índice temático}

% TODO (Daniel): completar cantidades y rangos; deben sumar 30 y cumplir
% 12 / 12 / 6 por dificultad.

{\scriptsize
\begin{longtable}{|p{6.5cm}|p{1.6cm}|p{1.6cm}|p{1.6cm}|p{1.6cm}|p{1.6cm}|}
    \hline
    \textbf{Sección del informe} & \textbf{Básicas} & \textbf{Interm.} & \textbf{Avanz.} & \textbf{Total} & \textbf{Rango} \\
    \hline
    \endhead
    1. Marco teórico y conceptual & & & & & \\ \hline
    2.1 Modelos de ejecución y aislamiento & & & & & \\ \hline
    2.2 Arranque en frío & & & & & \\ \hline
    2.3 Facturación por consumo y punto de equilibrio & & & & & \\ \hline
    2.4 Diseño orientado a eventos & & & & & \\ \hline
    2.5 WebAssembly y WASI & & & & & \\ \hline
    2.6 Límites técnicos y dependencia del proveedor & & & & & \\ \hline
    3. Análisis comparativo & & & & & \\ \hline
    4.1 Prueba de concepto & & & & & \\ \hline
    4.2 Modelo de costos & & & & & \\ \hline
    4.3 Caso TERABYTE & & & & & \\ \hline
    5. Tendencias y riesgos & & & & & \\ \hline
    \textbf{Total} & \textbf{12} & \textbf{12} & \textbf{6} & \textbf{30} & 1--30 \\ \hline
\end{longtable}
}

\subsection*{Preguntas}

%% Formato sugerido para cada pregunta (no hay macro en la clase):
%%
%% \textbf{P1 --- Básica --- Selección múltiple --- Sección 2.1.}
%% Enunciado. \\
%% a) ... \quad b) ... \quad c) ... \quad d) ... \\
%% \textit{Respuesta:} b. \textit{Justificación:} una o dos frases.
%%
%% Agrupar por sección de origen, en el mismo orden del índice temático.

% TODO (todos)

\clearpage

\end{document}
